\chapter{Results}

\section{Dataset \& Maneuver Statistics}
The dataset comprises 4,747 cyclists, from which 11,379 cyclist-cyclist interactions were identified after filtering for temporal overlap, proximity, and movement above walking speed. These interactions form the basis for maneuver extraction.

\subsection{Following Maneuvers}
\begin{center}
	\textbf{Missing Figure: Distribution plots of: duration, gap-min/-mean, thw min/mean}
\end{center}

From the filtered interactions, 3,792 following maneuvers were extracted, representing ~33\% of all cyclist-cyclist interactions. Following behaviors exhibit substantial variability: mean longitudinal gaps average 11.2m (mode ~5m, tail up to 35m), while minimum gaps often fall near 2m. Time headways are short (mean 2.7s, with many under 1s). Lateral offsets are tight (~0.63m), but speed changes range from -6m/s (deceleration) to +6m/s (acceleration), reflecting heterogeneous pacing and spacing strategies.

\subsection{Overtaking}
\begin{center}
	\textbf{Missing Figure: Distribution plots of: duration (pull out/merge back), lateral offset at cross, speed gain, speed mean}
\end{center}

A total of 573 overtaking maneuvers (~5\% of interactions) were detected. Most occur on the right side, with lateral offsets at crossing around 1.5m and maximum displacements exceeding 1.6m. Pull-out and merge-back durations average 3.8s, while speed gains vary broadly (mean 1m/s, extremes ±6m/s). Mean maneuver speeds are ~4.2m/s, highlighting a spectrum from cautious to aggressive behaviors.


\section{Stage 1: Behavioral Regimes}

\subsection{Overtaking}

\begin{center}
	\textbf{Missing Figure: Radar plot of cluster centers, both scenarios}
\end{center}

Overtaking maneuvers were segmented into 8,350 overlapping 1-second windows, of which 5,751 occurred at intersections and 2,599 on non-intersection segments. 
Although a grid search over candidate cluster numbers was conducted, the final selections were guided not solely by top-ranking metric values but by interpretability and the overall stability of clustering outcomes.
Clustering quality was high across both scenarios. For intersection windows, the silhouette score was 0.52 and the explained variance 0.75, with 91\% of windows assigned to the stable cluster and 9\% to the volatile cluster. Adjusted Rand Index metrics across subsamples and random seeds indicated consistent regime assignments (median ARI subsample = 0.84, ARI seed = 0.91), and median ARI under noise perturbation remained high (0.78). Non-intersection windows exhibited even stronger separation, with a silhouette score of 0.72 and explained variance of 0.72; 97\% of windows were stable and 3\% volatile. ARI values across subsamples and seeds were near-perfect (median ARI subsample = 1.00, ARI seed = 1.00), demonstrating extremely robust clustering.
¸
Feature selection identified a set of variables that contributed most meaningfully to scenario-specific clustering. 
Six features were common to both scenarios: \verb|acc_mad|, \verb|acc_max|, \verb|speed_qcv|, \verb|acc_qcv|, \verb|rot_fluc_max_abs| and \verb|rot_fluc_mad|. In addition, \verb|speed_mad| was retained only for intersection windows, whereas \verb|acc_min| and \verb|speed_std| were included exclusively for non-intersection windows.
Radar plots of cluster centers (Figure X) and correlation matrices (Appendix Table Y) summarize the feature-level distinctions.

The resulting clusters can be broadly characterized as stable and volatile, reflecting the magnitude and short-term variability of riding dynamics. The stable regime is marked by low fluctuations in speed, acceleration, and yaw rate, indicative of controlled, smooth riding behavior. In contrast, the volatile regime exhibits pronounced variability in these signals, reflecting abrupt or reactive maneuvering. Differences between scenarios are evident: intersection windows show greater variability across multiple features, consistent with higher environmental complexity and traffic interactions, while non-intersection windows are more homogeneous, reflecting smoother, less interrupted riding behavior. These findings demonstrate that overtaking maneuvers can be meaningfully decomposed into micro-behavioral regimes, with scenario-specific feature contributions and clearly distinguishable stability–volatility profiles.

\subsection{Following}
\begin{center}
	\textbf{Missing Figure: Radar plot of cluster centers, both scenarios}
\end{center}

In contrast to overtaking maneuvers, which exhibited clearly separated and event-like behavioral regimes, following behavior revealed a more gradual structure. While clustering outcomes were highly stable across resampling and initialization, separation between clusters was moderate, suggesting structured modulation of riding dynamics rather than sharply distinct behavioral modes.

Following maneuvers were segmented into 55,532 overlapping 3-second windows, with 26,297 windows (47.4\%) occurring at intersections and 29,235 (52.6\%) on non-intersection segments. Scenario-specific clustering was performed using riding-dynamics features selected through a metric-guided but interpretability-constrained procedure (see Methods). Configurations that maximized separation metrics occasionally produced degenerate solutions (e.g., clusters containing <1\% of windows) and were therefore discarded in favor of stable, non-degenerate partitions with meaningful size distributions.

At intersections, a three-cluster solution was selected. Clustering stability was exceptionally high, with median ARI subsample = 0.991 and ARI seed = 0.997, indicating near-identical partitions across subsamples and random initializations. Robustness to noise remained substantial (median ARI noise = 0.648). The explained variance was 0.762. Although the silhouette score (0.319) and Davies–Bouldin index (1.139) indicated moderate separation, cluster proportions were well balanced (48.5\%, 30.1\%, and 21.5\%), avoiding degeneracy and supporting structural validity.

For non-intersection segments, a two-cluster solution was selected. Stability metrics were again extremely high (median ARI subsample = 0.997, ARI seed = 0.9999), with ARI noise = 0.749. The explained variance reached 0.968, indicating that most variance in the feature space was captured by the partition. Separation metrics were lower (silhouette = 0.266; Davies–Bouldin = 1.668), but cluster sizes were nearly balanced (54.4\% and 45.6\%), suggesting persistent but moderately separated structure.

Across both scenarios, the retained features included \verb|acc_max|, \verb|acc_qcv|, \verb|speed_mad|, \verb|rot_fluc_max_abs| and \verb|rot_fluc_mad|. Intersection clustering additionally retained \verb|speed_qcv| and \verb|acc_mad|, whereas non-intersection clustering included \verb|acc_min| as a scenario-specific feature. Radar plots of cluster centers (Figure X) and full statistics (Appendix Table Y) illustrate feature-level distinctions.

Taken together, the combination of extremely high partition stability and only moderate geometric separation suggests that following behavior does not split into sharply distinct regimes, but rather into recurring patterns with gradual transitions between them. Such patterns can still be meaningfully represented as local states, as long as they reflect consistent and reproducible variations in observable riding dynamics. The resulting clusters therefore capture structured differences in motion variability without implying categorical behavioral types.

\paragraph{Intersection}
The largest cluster (48.5\%, R0) reflects \textbf{smooth, steady pedaling}, with low speed and acceleration variability, modest positive mean acceleration, and minimal short-term fluctuations—indicative of relaxed, consistent riding.

The second cluster (30.1\%, R1) represents \textbf{adjusting or flow riding}, with high relative acceleration variability and pronounced negative minima while mean acceleration is near zero. This captures frequent short-term braking and speeding up, without sustained directional change, as the rider modulates speed over terrain or traffic.

The third cluster (21.5\%, R2) corresponds to \textbf{sprint or burst efforts}, showing high speed dispersion and strong accelerative peaks, reflecting dynamic, high-intensity windows of rapid effort.

\paragraph{Non-intersection}

The largest cluster (54.4\%, R0) represents \textbf{dynamic cruising}, with moderate speed, higher acceleration variability, and noticeable short-term fluctuations. Riders exhibit frequent minor accelerations and decelerations, reflecting active longitudinal modulation while maintaining overall forward momentum.

The second cluster (45.6\%, R1) reflects \textbf{steady progression}, with slightly lower speed, reduced acceleration variability, and minimal short-term changes. This indicates smoother, more consistent riding with limited micro-adjustments.

These clusters are thus interpreted as differing primarily in intensity of longitudinal control variability, rather than categorical behavioral modes.
Compared to overtaking, following behavior exhibits higher regime multiplicity and lower geometric separation, suggesting a more continuous modulation of longitudinal control rather than event-driven volatility bursts.


\section{Aggregation Properties}

\begin{center}
	\textbf{Missing Figure: Plot of global-entropy, scenario exposures (both maneuvers), aggregated run lengths for all maneuvers}
\end{center}

\begin{center}
	\textbf{Missing Figure: Scatter plot: global entropy vs maneuver duration to show if structural diversity is just a duration artifact or real.}
\end{center}

Window-level regime assignments were aggregated to generate maneuver-level representations capturing the composition, persistence, and temporal distribution of local behavioral states. Aggregation considered scenario exposure, regime proportions, run-length statistics, centers of mass, and entropy measures (see Methods).

Following and overtaking maneuvers had comparable mean durations (7.6 s), though following maneuvers exhibited greater variability (SD = 8.0 s vs 2.2 s). Despite similar overall durations, the aggregation metrics reveal marked differences in structural composition.

\paragraph{Global Entropy and Diversity}
Global entropy was used to quantify the compositional diversity of regime sequences within individual maneuvers. While following and overtaking maneuvers exhibited comparable mean durations (7.6,s), their internal structure differed markedly.

Following maneuvers showed substantially higher heterogeneity (median entropy = 0.731; 38\% single-regime) compared to overtaking maneuvers (median entropy = 0.0; 54\% single-regime). Importantly, this difference cannot be attributed solely to maneuver length. Although entropy increased with duration in following maneuvers (Spearman $\rho$ = 0.54), overtaking maneuvers exhibited only a weak association ($\rho$ = 0.11). When restricting the analysis to the 5–10,s duration range, where both maneuver types are well represented, following maneuvers remained predominantly multi-regime (median entropy = 0.837), whereas overtaking maneuvers were still frequently dominated by a single regime (55\% single-regime).

These results indicate a clear difference in how maneuvers unfold over time. Following maneuvers typically consist of several shorter regime segments, reflecting ongoing adjustments during the interaction. Overtaking maneuvers, in contrast, are often dominated by a single regime that persists throughout most of the maneuver.

This pattern remains visible even when comparing maneuvers of similar duration, suggesting that the difference reflects distinct temporal organization rather than a simple length effect. In practical terms, overtaking tends to appear as a cohesive event, whereas following shows more internal variation.

These differences in internal structure motivate the maneuver-level analysis in Stage~2, where we examine whether such temporal patterns translate into distinct maneuver groups.

\paragraph{Scenario Exposure}
Overtaking maneuvers were largely concentrated in intersection segments (median intersection exposure = 1.0), while following maneuvers were more evenly distributed between intersection and non-intersection scenarios (median intersection exposure /approx 0.652). 
This pattern reflects that overtaking tends to occur as a localized, event-like maneuver, while following spans multiple segments. The observed distribution may also be influenced by differences in traffic density across the dataset, with intersections generally exhibiting higher traffic volumes than non-intersection segments.

\paragraph{Temporal Persistence and Center of Mass}
Run-length statistics reveal further structural differences between maneuver types. In overtaking maneuvers, persistent regimes frequently spanned the majority of the maneuver (mean run lengths up to /approx 9s), reflecting temporally uninterrupted local states. Following maneuvers, in contrast, displayed shorter and recurrent runs (mean run lengths /approx 2s), reflecting more fragmented and dynamically modulated state sequences.
For overtaking maneuvers, COMs show a notable spike around 0.5, consistent with maneuvers dominated by a single regime spanning the maneuver’s midpoint, alongside peaks at 0 for absent regimes. Following maneuvers display a less pronounced spike at 0.5, with a more broadly distributed COM around the middle of the maneuver. These patterns indicate that regimes can appear at varying temporal locations, and the COM distributions should be interpreted in the context of both regime presence and maneuver duration.


These structural differences motivate the second abstraction stage, which examines whether maneuvers cluster into distinct execution styles based on regime composition and temporal structure.

\section{Stage 2: Maneuver-Level Styles}
\subsection{Model Selection}
Maneuver-level styles were identified separately for following and overtaking maneuvers by clustering aggregated regime-based representations combined with maneuver-specific features. Feature selection and dimensionality reduction (PCA) were applied prior to K-means clustering, with hyperparameters chosen based on stability, as well as internal validity metrics and interpretability.


\subsection{Overtaking Maneuver Execution Styles}
For overtaking maneuvers, five clusters were obtained. Stability was near perfect: median ARI = 1.000 under subsampling and seed perturbations, and 0.993 under noise. The minimum cluster fraction was 0.027, indicating the smallest cluster was less frequent but still represented. Internal validity metrics were consistent with high-quality clustering: silhouette = 0.377, Calinski-Harabasz = 189, Davies-Bouldin = 1.02, and 60\% of variance explained across six principal components. Cluster sizes varied from 16 to 241 maneuvers, reflecting greater heterogeneity and imbalance compared to following maneuvers.


Overtaking maneuvers were clustered using aggregated regime composition, persistence, entropy, and maneuver-level dynamics. The five-cluster solution differentiates execution styles along three main dimensions: scenario exposure, micro-regime dominance, and maneuver intensity (speed, acceleration, lateral displacement).

\begin{center}
	\textbf{Missing Figure: Radar plot of cluster centers or heatmap, using most distinctive features}
\end{center}


Clustering reveals five distinct overtaking behavior types that differ primarily in speed level, acceleration strategy, lateral passing behavior, and temporal stability of behavioral regimes. These patterns broadly separate into slow intersection-related overtakes, fast free-flow overtakes, and rare high-steering maneuvers.

\paragraph{Slow and Moderate Intersection Overtakes}

Two clusters represent overtakes that occur predominantly at intersections, where cyclists typically move slower and interact more closely with other riders.

The first type (\textbf{Cluster 0}) consists of \textbf{slow and cautious overtakes}. These maneuvers are the shortest (duration $\approx 6.15\,s$) and occur at the lowest speeds ($speed_{mean} \approx 1.98\,m/s$), with minimal acceleration ($speed_{gain} \approx 0.27\,m/s$). The overtaken cyclist is also slow ($leader\_speed_{mean} \approx 1.31\,m/s$), and passing occurs at moderate distances ($distance_{cross} \approx 1.67\,m$). Regime dynamics show a balanced mix between behavioral states, indicating more frequent adjustments during the maneuver. Overall, this cluster represents careful passing in dense or constrained intersection environments.

The second intersection-related type (\textbf{Cluster 4}) represents \textbf{accelerating overtakes}. Starting from moderate speeds ($speed_{mean} \approx 3.90\,m/s$), cyclists accelerate strongly while passing ($speed_{gain} \approx 1.61\,m/s$; $long\_acc_{max} \approx 0.44\,m/s^2$, the highest among clusters). Maneuvers are slightly longer (duration $\approx 7.30\,s$) with balanced pull-out and merge-back phases. The passing gap is relatively wide ($lateral\_offset_{cross} \approx -0.59\,m$). Behavioral regimes remain highly stable throughout the maneuver, indicating a consistent overtaking strategy. This cluster therefore reflects controlled but decisive intersection overtakes where riders accelerate past slower cyclists.

\paragraph{Fast Overtakes in Mixed and Free-Flow Conditions}

Two clusters capture faster overtakes occurring in less constrained conditions, where cyclists maintain higher speeds and more stable motion.

The first of these (\textbf{Cluster 1}) represents \textbf{fast and wide overtakes} with sustained effort. These maneuvers are the longest (duration $\approx 8.91\,s$) and occur at high speeds ($speed_{mean} \approx 5.30\,m/s$; $speed_{max} \approx 6.03\,m/s$). The overtaker maintains a clear speed advantage over the cyclist being passed ($rel\_speed_{mean} \approx 1.70\,m/s$). Riders also move further laterally when overtaking ($lateral\_offset_{cross} \approx -0.53\,m$), resulting in relatively long overtaking distances ($overtake\_length \approx 14.79\,m$). Behavioral regimes are extremely stable, suggesting a steady, uninterrupted overtaking strategy.

The second type (\textbf{Cluster 3}) represents fast but \textbf{smooth free-flow overtakes}. This cluster occurs almost exclusively outside intersections and shows the highest average speed ($speed_{mean} \approx 5.35\,m/s$) while maintaining very low variability ($speed_{std} \approx 0.35\,m/s$). The overtaken cyclist is also relatively fast ($leader\_speed_{mean} \approx 3.82\,m/s$), and passing occurs with a consistent speed advantage ($rel\_speed_{mean} \approx 1.91\,m/s$). Compared to Cluster 1, lateral displacement is smaller ($lateral\_offset_{cross} \approx -0.30\,m$) and the maneuver appears more streamlined. Overall, these overtakes represent smooth passing in free-flow cycling conditions where both riders travel quickly.

\paragraph{High-Steering Overtakes}

Finally, a small but distinctive cluster (\textbf{Cluster 2}) captures overtakes characterized by \textbf{sharp steering}. These maneuvers are relatively short (duration $\approx 5.84\,s$) and occur at moderately high speeds ($speed_{mean} \approx 4.80\,m/s$), but they stand out due to extremely high lateral movement ($lat\_speed_{max} \approx 2.23\,m/s$; $lat\_acc_{max} \approx 0.115\,m/s^2$), far exceeding values in other clusters. The overtaken cyclist travels at a similar speed ($leader\_speed_{mean} \approx 4.41\,m/s$), resulting in only a modest speed advantage. Behavioral regimes alternate more frequently than in fast free-flow overtakes. This cluster likely reflects situations where cyclists must steer sharply to initiate the pass, for example when space is limited or when overtaking riders with similar speeds.

\paragraph{Summary}

Taken together, the clusters highlight three overarching overtaking patterns:
\begin{itemize}
	\item Slow intersection passes with cautious or accelerating behavior (Clusters 0 and 4),
	\item Fast, stable overtakes in mixed or free-flow conditions (Clusters 1 and 3),
	\item Rare steering-intensive overtakes characterized by abrupt lateral movement (Cluster 2).
\end{itemize}

These patterns suggest that overtaking behavior is strongly shaped by both traffic context and relative cyclist speeds, which in turn influence how riders manage acceleration, lateral clearance, and maneuver duration.

\subsection{Following Maneuver Execution Styles}

For following maneuvers, six clusters were identified. Stability measures indicated robust separation: median subsample ARI = 0.933, seed ARI = 0.987, and noise ARI = 0.923. The minimum cluster fraction was 0.111, suggesting no cluster was sparsely populated. Internal validity metrics supported the clustering choice, with a silhouette score of 0.245, Calinski-Harabasz index of 616, Davies-Bouldin index of 1.40, and 61\% of variance explained by the first seven principal components. Cluster sizes ranged from 430 to 757 maneuvers, reflecting relatively balanced partitions.

\begin{center}
	\textbf{Missing Figure: Radar plot of cluster centers or heatmap, using most distinctive features}
\end{center}


paragraph{Short and Moderate Following}

Two clusters represent following behavior over short to moderate durations, with cyclists maintaining moderate gaps and adjusting their speed to track the leader.

The first type (\textbf{Cluster 4}) corresponds to brief, stable following. These maneuvers are short (duration $\approx 4.84,s$) with moderate spacing ($gap_{mean} \approx 11.18,m$) and mild acceleration ($speed_{gain} \approx 0.32,m/s$). Cyclists maintain a consistent position relative to the leader, with moderate speed correlation ($speed_{corr} \approx 0.33$), indicating balanced responsiveness. This cluster reflects cautious, short interactions where the follower maintains a steady gap without large speed adjustments.

The second type (\textbf{Cluster 0}) represents slightly longer but still moderate following. Cyclists maintain larger gaps ($gap_{mean} \approx 12.49,m$) and show more responsive speed adjustments ($speed_{gain} \approx 0.69,m/s$; $speed_{corr} \approx 0.61$), allowing smoother adaptation to the leader’s motion. These maneuvers last slightly less than 5 s on average (duration $\approx 4.61,s$) and involve low lateral variability ($lateral_offset_{std} \approx 0.52,m$). This cluster captures moderate tracking behavior with a balance between safety and responsiveness.

\paragraph{Quick and Catch-Up Following}

Two clusters capture maneuvers with higher acceleration and tighter spacing, where followers actively adjust to close gaps.

The first of these (\textbf{Cluster 2}) represents short, reactive adjustments. Maneuvers are brief (duration $\approx 6.17,s$) with close spacing ($gap_{mean} \approx 8.25,m$) and strong acceleration ($speed_{gain} \approx 1.20,m/s$). Followers respond quickly to the leader’s motion (speed correlation $speed_{corr} \approx 0.82$; response delay $response_delay \approx 0.44,s$), allowing rapid catch-up. This cluster reflects high-intensity, short-term adjustments to maintain or reduce headway.

The second type (\textbf{Cluster 1}) captures sustained catch-up behavior. These maneuvers are longer (duration $\approx 7.43,s$) with relatively low minimum headway ($thw_{min} \approx 1.85,s$) and positive relative speed ($rel_speed_{mean} \approx 0.21,m/s$), indicating the follower accelerates above the leader’s speed. Maximum speeds are high ($speed_{max} \approx 5.21,m/s$), and lateral variation is moderate ($lateral_offset_{std} \approx 0.35,m$). This cluster likely reflects followers actively closing distance over medium durations, e.g., after temporary slowing or in dense traffic.

\paragraph{Prolonged and Spaced Following}

Two clusters reflect longer, steady tracking, with either close or wide spacing.

The first (\textbf{Cluster 3}) represents prolonged, close following. Maneuvers are the longest (duration $\approx 20.95,s$) with the tightest spacing ($gap_{min} \approx 6.08,m$) and relatively slow follower response ($response_delay \approx 1.27,s$). Acceleration is mild ($acc_{mean} \approx 0.01,m/s^2$), indicating steady maintenance rather than active catch-up. This cluster characterizes sustained tracking behind a leader over long stretches, where lateral adjustments are moderate ($lateral_offset_{std} \approx 0.42,m$).

The second (\textbf{Cluster 5}) corresponds to spaced cruising. These maneuvers have moderately long duration (duration $\approx 7.01,s$) and the largest gaps ($gap_{mean} \approx 13.70,m$). Speed variability is low ($speed_{std} \approx 0.20,m/s$) and coupling with leader speed is weak ($speed_{corr} \approx 0.21$), suggesting relaxed, low-interaction following. This cluster reflects situations where cyclists maintain wide separation and stable speed, for example in lightly trafficked sections or relaxed group riding.




\subsection{Interpretation of Maneuver-Level Clusters}
Taken together, these maneuver-level clusters translate the micro-regime structure identified in Stage 1 into coherent patterns of execution. While the clusters largely reflect scenario exposure and dominant regime composition, they provide a clear picture of how maneuvers are organized over time: some are almost entirely composed of a single stable or volatile regime, others alternate between regimes, and a few span mixed contexts. These patterns confirm that micro-regimes are not isolated fluctuations but tend to form persistent sequences at the maneuver level.

The value of this stage lies in summarizing maneuver structure in terms of persistence, variability, and context, offering a higher-level view that complements the detailed Stage 1 characterization. These execution-style clusters provide a practical framework for grouping and comparing maneuvers, quantifying temporal fragmentation, and linking behavior to environmental conditions, without implying fundamentally new behavioral modes.

To provide an overview of the identified maneuver types, Table X summarizes cluster sizes, contextual occurrence probabilities, dominant behavioral regimes, and the main behavioral characteristics distinguishing the clusters.


% Please add the following required packages to your document preamble:
% \usepackage{multirow}
% \usepackage{graphicx}
\begin{table}[]
	\caption{Summary of identified overtaking and following maneuver clusters, with their prevalence, regime composition, and key behavioral characteristics.}
	\resizebox{\textwidth}{!}{%
		\begin{tabular}{|c|c|c|l|c|c|l|m{4cm}|}
			\hline
			$\tau$                      & Cluster & count & Short Label      & $P_m(I)$ & $P_m(N)$ & Regime composition                         & Behavioral characteristics                                      \\ \hline
			\multirow{5}{*}{Overtaking} & 0       & 87    & Slow Steady      & 0.98     & 0.02     & stable 48,3\% , volatile 51,7\%            & Gradual, low-speed, evenly paced maneuver                       \\ \cline{2-8} 
			& 4       & 241   & Accelerating     & 0.98     & 0.02     & stable 98,9\%                              & Moderate-speed with steady longitudinal acceleration            \\ \cline{2-8} 
			& 1       & 141   & Fast Wide        & 0.56     & 0.44     & stable (I, 99,1\%),  stable(N, 99,6\%)     & High-speed, wide lateral displacement, sustained effort         \\ \cline{2-8} 
			& 3       & 88    & Fast Smooth      & 0.00     & 1.00     & stable 99,9\%                              & High-speed, steady, streamlined, minimal lateral variation      \\ \cline{2-8} 
			& 2       & 16    & Sharp Steering   & 0.03     & 0.97     & stable 47\%, volatile 53\%                 & Abrupt, lateral-dominated, agile maneuvering                    \\ \hline
			\multirow{6}{*}{Following}  & 5       & 635   & Spaced Cruise    & 0.04     & 0.96     & N-R0 80\%, N-R1  20\%                      & Moderate duration, largest gaps, stable low-variability speeds  \\ \cline{2-8} 
			& 4       & 721   & Short Steady     & 0.88     & 0.12     & I-R0 15\%, I-R1 80\%                       & Short-duration, moderate gap, mild acceleration                 \\ \cline{2-8} 
			& 0       & 757   & Moderate         & 0.88     & 0.12     & I-R2 91\%                                  & Medium gaps, mild speed adjustments, balanced responsiveness    \\ \cline{2-8} 
			& 3       & 430   & Prolonged Close  & 0.48     & 0.52     & I-R0 50\%, I-R1 45\%, N-R0 56\%, N-R1 37\% & Long-duration, low acceleration, close spacing                  \\ \cline{2-8} 
			& 1       & 560   & Catch-Up         & 0.25     & 0.75     & I-R0 37\%, I-R1 10\%, N-R0 15\%, N-R1 85\% & Follower accelerates to reduce gap, short THW, higher max speed \\ \cline{2-8} 
			& 2       & 689   & Quick Adjustment & 0.89     & 0.11     & I-R0 88\%                                  & Short, tight spacing, rapid acceleration, fast response         \\ \hline
		\end{tabular}%
	}
\end{table}

\section{Infrastructure \& Traffic Conditioning}

Infrastructure characteristics showed little variation across maneuver clusters within the same scenario category. This is expected because windows were already separated into intersection and non-intersection contexts during regime discovery. Consequently, infrastructure variables were largely homogeneous within each scenario group and did not contribute strongly to maneuver-level differentiation.

Due to the small sample size (n=16), ECDF curves exhibit strong step artifacts and no robust contextual interpretation can be derived.
\iffalse
found a few ecdf plots with distinct cluster curves: overtaking: for intersection dominant clusters: pedestrian_drac: cluster 0 rises fastest with low drac, followed by cluster 4 and then cluster 1 bicycle_drac: cluster 0 rises fastest again, followed by cluster 1 and then 0 bicycle count (within 5m): cluster 1 and 4 experience higher counts compared to 0 car and pedestrian min ttc: cluster 4 has rises faster with low min-ttc compared to 0 and 1 non-intersection dominant: (including cluster 1 since mixed contexts) bicycle count: cluster 1 experiences higher bicycle counts compared to 3 and 2 stepwise nature of cluster 2 (16 maneuvers) makes it hard to interpret most, (does it lay above or below) following: when comparing all together 1 diverges most, but probably since it is has largest scenario spread, and it follows a different curve than the other clusters intersection-dominant: for clusters 0,2,4 most curves look identical only slight differences how would i show this? i cant just take some random plot and say "look here, no difference" non-intersection-dominant: (3 und 5) bicycle-min-ttc: cluster 5 experiences a lot lower min-ttc than 3 fraction-of-ttc-below-3s-pedestrian: both have an s curve but teh s-curve for 5 is shallow while the one for 3 is steep. -> 3 experiences more higher fractions of ttcs below th min-ttc car and pedestrian: cluster 5 sooner reaches the right border (max ttc of 6s, higher ttc are not considered) meaning 5 experiences less min-ttc between 2 and 6 seconds
\fi